\section{Mechanisms and Alternative Explanations}
\label{sec:mechanisms}

We examine three mechanisms to distinguish the cover bidding hypothesis
from alternative explanations.

\subsection{M1: Strategic Selection vs.\ Competitive Displacement}

A key concern is whether FL presence \emph{causes} higher prices or
merely correlates with tender characteristics that independently drive
prices up. Under competitive displacement, FL firms would crowd out
genuine competitors, reducing effective competition. Under strategic
selection, cartels deploy cover bidders in tenders that already attract
genuine competition, and the cover bidders are added on top.

We regress the number of genuine (non-FL) competitors on FL presence.
The OLS coefficient is 0.143 ($p < 0.01$): FL-present tenders have
\emph{more}, not fewer, genuine competitors. This rules out competitive
displacement.

The IV estimate for non-FL firms (Table~\ref{tab:iv_main}, column~4) is
0.404 ($p < 0.01$), consistent with the OLS finding: instrumented FL presence
is associated with substantially more genuine competitors. Cartels
appear to select into competitive tenders and add cover bidders on top
of existing competition. This rules out competitive crowding-out and is
consistent with FL bidders being added on top of genuine competition, as
cover-bidding theory predicts.

\subsection{M2: Reference Price Calibration}

If cartels use cover bids to inflate reference prices for future
tenders, winners in FL-present tenders should bid closer to the
reference price. We regress $\log(p_{\text{negot}} / p_{\text{ref}})$ on
FL presence. The coefficient is $-0.041$ ($p < 0.01$), meaning
that winning bids in FL-present tenders are 4\% closer to the reference
price. This is consistent with cartels calibrating the winning bid near
the reference price while cover bidders submit bids that do not
undercut it. However, this pattern could also reflect item % REVISION-V5: added item-composition caveat per referee minor 8
composition: FL-present tenders may involve more standardized goods
whose prices are naturally closer to the reference price. We cannot
fully rule out this alternative without controlling for item-level
price-to-reference dispersion, which our data do not support at
sufficient granularity.

\subsection{M3: Reverse Causality}

Could FL firms be attracted to high-priced tenders rather than causing
higher prices? We test this by regressing FL entry on \emph{lagged} log
prices within market (item $\times$ PBU) using a panel of 427,474
market-year observations. The coefficient is 0.0021 (SE~=~0.0008,
$p = 0.012$)---statistically significant but economically negligible.
In elasticity terms, a 10\% increase in lagged prices is associated with
a $0.1 \times 0.0021 = 0.00021$ increase in FL entry probability, or
0.02 percentage points against an unconditional FL rate of 4.8\%---an
elasticity of $0.00021 / 0.048 \approx 0.004$, essentially zero. This
small magnitude suggests that reverse causality---while detectable in a
sample of 1.6 million observations---cannot explain the 6--9\% price
effect attributed to FL presence. The reverse-causality effect is two
orders of magnitude smaller than the main estimate, indicating that
price-setting selection into FL status does not drive our results.
Table~\ref{tab:mechanisms} summarizes all mechanism tests.

\begin{table}[htbp]
\centering
\caption{Mechanism Tests}
\label{tab:mechanisms}
\begin{adjustbox}{max width=\textwidth}
\begin{threeparttable}
\small
\begin{tabular}{lccc}
\toprule
Mechanism & Coefficient & SE & N \\
\midrule
M1: Competitive displacement (OLS) & 0.1426*** & (0.0062) & 1,654,447 \\
\quad DV: $\log(n_{\text{genuine}} + 1)$ & & & \\
M1: Competitive displacement (IV) & 0.3244*** & (0.0249) & 1,654,447 \\
M2: Reference price calibration & -0.0412*** & (0.0031) & 1,654,394 \\
\quad DV: $\log(p_{\text{negot}} / p_{\text{ref}})$ & & & \\
M3: Reverse causality & 0.0021** & (0.0008) & 427,474 \\
\bottomrule
\end{tabular}
\begin{tablenotes}
\small
\item \textit{Notes:} M1: FL presence on genuine competition. M2: bid-to-reference ratio.
M3: lagged price predicting FL entry (reverse causality check).
*** \textit{p}$<$0.01, ** \textit{p}$<$0.05, * \textit{p}$<$0.1.
\end{tablenotes}
\end{threeparttable}
\end{adjustbox}
\end{table}


\subsection{Interpretation of the Convite--Preg\~ao Difference}

The FL coefficient is larger for preg\~ao (9.3\%) than for convite
(3.8\%). This may appear counterintuitive if preg\~ao's electronic
format provides more transparency. However, the difference reflects
two institutional features. First, preg\~ao allows real-time observation
of competing bids, which facilitates Regime~1 cover bidding: cover
bidders can monitor the auction and ensure their bids remain above the
designated winner's bid. Second, convite's sealed-bid format makes it
harder for cover bidders to calibrate bids---but also creates less need
for cover bidding when there are already three genuine bidders (the
minimum threshold). The smaller convite coefficient thus reflects both
reduced scope for cover bidding under sealed-bid format and the
compositional effect of the minimum-bidder rule.

\subsection{Alternative Explanations} % REVISION-V5: new section per referee 5.5

We consider four non-collusive explanations for FL status and test each
against observable data.

\emph{Financial distress.} Firms near bankruptcy may bid on many tenders
without winning, generating FL status without collusion. Using firm size
(porte) from the BEC registry as a proxy, we find that FL firms are
slightly more likely to be micro or small enterprises (77.2\% vs.\
73.0\% for non-FL always-losers, $p < 0.001$). However, the difference
is modest and the FL--price effect persists after controlling for
Imhof-style bid-level screen indicators (Section~\ref{sec:robustness}),
suggesting that firm-size composition alone does not explain the price
anomaly.

\emph{Learning-by-doing.} New entrants may participate repeatedly while
learning the procurement system. Using the firm registry's activity start
date, we find that FL firms are on average \emph{younger} than non-FL
always-losers (14.6 vs.\ 15.6 years, $p < 0.001$), which is
directionally consistent with a learning channel. However, the magnitude
is small (one year) and the combined logit controlling for age,
geography, and size shows that age explains only a small fraction of FL
status ($z = -3.5$, marginal effect $\approx -0.008$ per year of age).

\emph{Geographic distance.} FL firms that are geographically distant
from PBUs may consistently lose to local competitors. Among FL firms,
76.8\% are headquartered in S\~ao Paulo state, compared to 82.2\% of
non-FL always-losers ($p < 0.001$). FL firms are modestly more likely
to be out-of-state, consistent with some geographic disadvantage, though
the 5.4 percentage point difference is insufficient to explain the
6--9\% price effect.

\emph{Minimum-bidder compliance.} PBU officials may invite firms to
meet convite's three-bidder minimum, creating FL status through
institutional compliance rather than collusion. This channel is
plausible for convite but cannot explain the FL--price effect in
preg\~ao (which has no minimum-bidder rule), where the coefficient is
actually larger (9.3\% vs.\ 3.8\%).

Table~\ref{tab:alternative_mechanisms} in
Appendix~\ref{app:heterogeneity} reports the full results.

\subsection{Summary}

The mechanism tests point consistently toward FL presence predicting
price inflation through a process that is difficult to explain by
competitive behavior alone. No single mechanism is conclusive in
isolation, but their joint convergence---together with the ruling-out of
alternative explanations---substantially narrows the space of non-collusive
interpretations, even as we acknowledge that definitive causal
attribution remains elusive.
